\section{What can be used, and what needs a different estimator?}
The derivative direction comes from independent directional differences,
not from the iso score. Its error mechanisms are concrete: curvature bias
decreases on smaller well-shaped cells, while noise is amplified by the
inverse edge matrix. Conditioning and noise scale can therefore inform a
search decision.

Recursive lifting is an understandable research construction, but every
stage changes the field and support. Raw mixed partials, symmetry defect
and Laplacian are useful outputs to inspect. The clean quadratic counterexample
prevents treating them as validated curvature. Larger local polynomial
neighborhoods, recovered gradient regularization or a potential-based fit
are appropriate next investigations when curvature accuracy matters.

Iso coverage visualizes where selected derivative levels cross the mesh.
Its counting measure, level selection and bandwidth are part of its
definition. The affine counterexample and contamination tests show why
this descriptor cannot generally serve as confidence. Reversing the anomaly
sign on one dataset would merely substitute another uncalibrated rule.

A measurement system should combine a declared metric, derivative estimates,
support and noise diagnostics, application constraints, and acceptance
against a new objective measurement. Confidence additionally needs held-out
accuracy targets, independent repetitions or known measurement uncertainty,
and validation across curvature, support and contamination regimes.
The Python package makes these distinctions reproducible; it does not yet
provide an experimentally validated controller.
